\documentclass[11pt,a4paper]{article}

% ── Packages ──
\usepackage[utf8]{inputenc}
\usepackage[T1]{fontenc}
\usepackage{geometry}
\usepackage{hyperref}
\usepackage{xcolor}
\usepackage{listings}
\usepackage{tcolorbox}
\usepackage{fancyhdr}
\usepackage{enumitem}
\usepackage{booktabs}
\usepackage{titlesec}
\usepackage{parskip}

% ── Page Setup ──
\geometry{margin=1in}
\hypersetup{
    colorlinks=true,
    linkcolor=blue,
    urlcolor=cyan,
    citecolor=blue
}

% ── Header/Footer ──
\pagestyle{fancy}
\fancyhf{}
\fancyhead[L]{\textbf{Gradescope AutoGrade}}
\fancyhead[R]{User Manual v0.1.0}
\fancyfoot[C]{\thepage}
\renewcommand{\headrulewidth}{0.5pt}

% ── Code Style ──
\lstdefinestyle{mystyle}{
    backgroundcolor=\color{gray!5},
    basicstyle=\ttfamily\small,
    breaklines=true,
    frame=single,
    rulecolor=\color{gray!30},
    numbers=none,
    xleftmargin=1em,
    framexleftmargin=1em,
}
\lstset{style=mystyle}

\lstdefinestyle{yamlstyle}{
    backgroundcolor=\color{yellow!5},
    basicstyle=\ttfamily\small,
    breaklines=true,
    frame=single,
    rulecolor=\color{yellow!50!black},
    numbers=none,
    xleftmargin=1em,
    framexleftmargin=1em,
}

% ── Custom Boxes ──
\newtcolorbox{tipbox}[1][]{
    colback=green!5,
    colframe=green!50!black,
    title=\textbf{Tip},
    #1
}

\newtcolorbox{warnbox}[1][]{
    colback=red!5,
    colframe=red!50!black,
    title=\textbf{Warning},
    #1
}

\newtcolorbox{infobox}[1][]{
    colback=blue!5,
    colframe=blue!50!black,
    title=\textbf{Info},
    #1
}

% ── Section Title Format ──
\titleformat{\section}{\Large\bfseries}{\thesection}{1em}{}
\titleformat{\subsection}{\large\bfseries}{\thesubsection}{1em}{}

\begin{document}

% ── Title Page ──
\begin{titlepage}
\centering
\vspace*{3cm}

{\Huge\bfseries Gradescope AutoGrade\par}
\vspace{0.5cm}
{\Large User Manual\par}
\vspace{0.3cm}
{\large Version 0.1.0\par}

\vspace{2cm}

{\large AI-powered automated grading assistant for TAs\par}

\vspace{3cm}

\begin{tabular}{rl}
\textbf{Repository:} & \url{https://github.com/tzhazuma/gradescope_autograde} \\
\textbf{License:} & GPL-3.0 \\
\textbf{Platform:} & macOS (primary), Linux (compatible) \\
\textbf{Python:} & 3.11+ (3.12 recommended) \\
\end{tabular}

\vfill
{\large June 2026\par}
\end{titlepage}

% ── Table of Contents ──
\tableofcontents
\newpage

% ════════════════════════════════════════════════════════════
\section{Introduction}

Gradescope AutoGrade is an AI-powered tool that automates the grading workflow on Gradescope. It fetches student submissions, evaluates them using large language models (LLMs) against configurable rubrics, and uploads scores back to Gradescope — all with minimal human intervention.

\subsection{Who Is This For?}

This tool is designed for \textbf{Teaching Assistants (TAs)} and \textbf{instructors} who:
\begin{itemize}
    \item Regularly grade large numbers of student submissions on Gradescope
    \item Want to reduce repetitive grading work
    \item Need consistent, rubric-based evaluation
    \item Prefer to keep a human in the loop for uncertain cases
\end{itemize}

\subsection{Key Features}

\begin{itemize}[itemsep=3pt]
    \item \textbf{Automated submission fetching} — downloads student answers from Gradescope via HTTP scraping
    \item \textbf{AI grading} — evaluates answers against YAML-defined rubrics using LLMs
    \item \textbf{Multiple LLM backends} — OpenCode Go (cloud), LM Studio (local/private), OpenAI, Anthropic
    \item \textbf{PDF question parsing} — extracts questions and reference answers from instructor PDFs
    \item \textbf{Confidence-based review queue} — flags uncertain grades (confidence $<$ 70\%) for human review
    \item \textbf{Progress tracking} — resumes interrupted grading runs without re-processing
    \item \textbf{CSV/JSON export} — exports results in Gradescope-compatible format
    \item \textbf{Interactive TUI \& Web GUI} — Terminal UI and cross-platform web interface
    \item \textbf{LM Studio auto-management} — Auto-detect, start, and stop local models
    \item \textbf{10 CLI commands} — full workflow via terminal
\end{itemize}

\subsection{How It Works}

\begin{infobox}
\textbf{Important:} Gradescope does \textbf{not} provide a public API. All interaction is performed by scraping the same HTML pages you see in your browser, using your instructor/TA credentials.
\end{infobox}

The tool operates in four layers:

\begin{enumerate}[itemsep=3pt]
    \item \textbf{Transport Layer} — Manages HTTP sessions, CSRF authentication, cookie persistence, and rate limiting.
    \item \textbf{Client Layer} — Scrapes Gradescope HTML pages to list courses, assignments, and submissions, and to upload grades.
    \item \textbf{Grading Engine} — Sends rubrics and student answers to an LLM, parses the structured JSON response.
    \item \textbf{Workflow Orchestrator} — Coordinates the full pipeline: fetch $\rightarrow$ grade $\rightarrow$ review $\rightarrow$ export $\rightarrow$ upload.
\end{enumerate}

% ════════════════════════════════════════════════════════════
\section{Installation}

\subsection{Prerequisites}

\begin{itemize}[itemsep=2pt]
    \item \textbf{Python 3.11+} (Python 3.12 is recommended)
    \item \textbf{macOS} (primary platform; Linux is largely compatible)
    \item \textbf{Gradescope account} with instructor or TA access
\end{itemize}

\subsection{Setup Steps}

\begin{enumerate}[itemsep=5pt]
    \item Clone the repository:
\begin{lstlisting}[language=bash]
git clone https://github.com/tzhazuma/gradescope_autograde.git
cd gradescope_autograde
\end{lstlisting}

    \item Create and activate a virtual environment:
\begin{lstlisting}[language=bash]
python3.12 -m venv venv
source venv/bin/activate
\end{lstlisting}

    \item Install the package:
\begin{lstlisting}[language=bash]
pip install -e .
\end{lstlisting}

    \item (Optional) Install browser automation support:
\begin{lstlisting}[language=bash]
pip install -e ".[browser]"
playwright install chromium
\end{lstlisting}
\end{enumerate}

\subsection{Verify Installation}

\begin{lstlisting}[language=bash]
gs-autograde --help
\end{lstlisting}

You should see a list of 10 available commands: \texttt{export}, \texttt{grade}, \texttt{gui}, \texttt{list}, \texttt{login}, \texttt{models}, \texttt{parse-pdf}, \texttt{review}, \texttt{tui}, \texttt{upload}.

% ════════════════════════════════════════════════════════════
\section{Configuration}

\subsection{Creating Your Config}

Copy the example configuration and edit it:

\begin{lstlisting}[language=bash]
cp config/config.example.yaml config/config.yaml
\end{lstlisting}

\begin{warnbox}
\texttt{config/config.yaml} is git-ignored. Never commit your credentials to version control.
\end{warnbox}

\subsection{Minimal Configuration}

For a quick start, you only need to set your email and password:

\begin{lstlisting}[style=yamlstyle]
auth:
  email: "your-email@university.edu"
  password: "your-password"

llm:
  provider: "opencode-go"
  model: "deepseek-v4-flash"
\end{lstlisting}

Set your API key:

\begin{lstlisting}[language=bash]
export OPENCODE_GO_API_KEY="your-api-key"
\end{lstlisting}

\subsection{Full Configuration Reference}

\begin{lstlisting}[style=yamlstyle]
# ── Authentication ──
auth:
  email: "your-email@university.edu"
  password: ""
  cookie: ""                          # Browser session cookie

# ── Gradescope Settings ──
gradescope:
  base_url: "https://www.gradescope.com"
  request_delay: 1.0                  # Seconds between requests
  max_retries: 3

# ── LLM Grading Engine ──
llm:
  provider: "opencode-go"             # opencode-go|lmstudio|openai|anthropic
  model: "deepseek-v4-flash"
  multimodal_model: "mimo-v2.5"
  api_key: "${OPENCODE_GO_API_KEY}"
  base_url: "https://opencode.ai/zen/go/v1"
  temperature: 0.1
  max_tokens: 4096

# ── LM Studio (Local) ──
lmstudio:
  base_url: "http://localhost:1234/v1"
  auto_discover: true
  default_model: ""

# ── Workflow ──
workflow:
  auto_upload: false
  review_threshold: 0.7               # Confidence < 0.7 → review
  batch_size: 50

# ── Output ──
output:
  grades_dir: "data/output/grades/"
  format: "gradescope_csv"
  generate_feedback: true
\end{lstlisting}

% ════════════════════════════════════════════════════════════
\section{CLI Commands}

\subsection{\texttt{login} — Authenticate}

\begin{lstlisting}[language=bash]
# Interactive login (prompts for credentials)
gs-autograde login

# Email-based login
gs-autograde login --email user@university.edu

# Cookie-based (from browser DevTools → Application → Cookies)
gs-autograde login --cookie "_gradescope_session=abc123..."

# Browser-based (opens Chromium, you log in manually, cookies auto-captured)
gs-autograde login --browser

# Browser-based with custom timeout (default: 120s)
gs-autograde login --browser --timeout 300
\end{lstlisting}

Successful login saves cookies to \texttt{.cookies/session.txt}.

The \texttt{--browser} flag opens a Chromium window (via Playwright) where you log into Gradescope manually. After successful login, the tool automatically extracts session cookies from the browser. This is the most reliable login method as it handles any CAPTCHA or 2FA challenges that Gradescope may present. Requires Playwright: \texttt{pip install -e ".[browser]"}.

\subsection{\texttt{list} — Browse Resources}

\begin{lstlisting}[language=bash]
gs-autograde list courses
gs-autograde list assignments <COURSE_ID>
gs-autograde list submissions <COURSE_ID> <ASSIGNMENT_ID>
\end{lstlisting}

Outputs results as formatted Rich tables.

\subsection{\texttt{models} — Discover AI Models}

\begin{lstlisting}[language=bash]
gs-autograde models
\end{lstlisting}

Displays a table of available models from all configured providers:

\begin{center}
\begin{tabular}{llrlr}
\toprule
\textbf{Provider} & \textbf{Model} & \textbf{Context} & \textbf{Multimodal} & \textbf{Loaded} \\
\midrule
opencode-go & DeepSeek V4 Flash & 1,000,000 & No & $\checkmark$ \\
opencode-go & MiMo V2.5 & 1,000,000 & Yes & $\checkmark$ \\
lmstudio & Gemma 4 12B & 131,072 & Yes & $\checkmark$ \\
lmstudio & Qwen 3.5 9B & 262,144 & No & $\checkmark$ \\
\bottomrule
\end{tabular}
\end{center}

For LM Studio, models are queried in real-time via \texttt{GET /api/v1/models}. Only downloaded models are shown, with a checkmark indicating whether they are loaded in memory.

\subsection{\texttt{parse-pdf} — Extract Questions}

\begin{lstlisting}[language=bash]
# Display extracted questions
gs-autograde parse-pdf exam_questions.pdf

# Save as rubric template
gs-autograde parse-pdf exam_questions.pdf -o rubric.yaml

# Custom question separator
gs-autograde parse-pdf exam_questions.pdf --separator "### Problem"
\end{lstlisting}

\begin{infobox}
The PDF should separate questions with a consistent marker like \texttt{\#\# Question}. The answer section begins after a blank line or an \texttt{Answer:} prefix. Edit the generated YAML to add per-question rubric criteria and point values.
\end{infobox}

\subsection{\texttt{grade} — Run Grading Pipeline}

This is the core command. It fetches submissions, grades them, and optionally uploads scores.

\begin{lstlisting}[language=bash]
# Dry run: grade locally, do NOT upload
gs-autograde grade 123456 789012 \
  --rubric rubric.yaml \
  --dry-run

# Full run with OpenCode Go
gs-autograde grade 123456 789012 \
  --rubric rubric.yaml \
  --provider opencode-go \
  --model deepseek-v4-flash

# Full run with local LM Studio model
gs-autograde grade 123456 789012 \
  --rubric rubric.yaml \
  --provider lmstudio \
  --model gemma4-12b
\end{lstlisting}

\begin{tipbox}
\textbf{Always start with \texttt{--dry-run}} to verify your rubric produces reasonable grades before uploading anything to Gradescope.
\end{tipbox}

\subsection{\texttt{review} — Review Flagged Grades}

\begin{lstlisting}[language=bash]
gs-autograde review
\end{lstlisting}

Presents each flagged submission interactively:
\begin{itemize}
    \item Student name and question
    \item AI-assigned score with confidence level
    \item AI-generated feedback
    \item Options: \texttt{[a]pprove}, \texttt{[r]eject} (with notes), \texttt{[s]kip}
\end{itemize}

\subsection{\texttt{upload} — Submit Grades}

\begin{lstlisting}[language=bash]
gs-autograde upload 123456 789012 --results results.json
\end{lstlisting}

\subsection{\texttt{export} — Export Results}

\begin{lstlisting}[language=bash]
gs-autograde export --format gradescope    # Gradescope bulk upload CSV
gs-autograde export --format detailed      # Per-criterion breakdown
gs-autograde export --format json          # Full JSON
\end{lstlisting}

\subsection{\texttt{tui} --- Interactive Terminal UI}

\begin{lstlisting}[language=bash]
gs-autograde tui
\end{lstlisting}

Launches a full-screen Textual-based terminal interface. The TUI guides you through:

\begin{enumerate}[itemsep=3pt]
    \item \textbf{Course Selection} --- Browse and select from your Gradescope courses using keyboard navigation.
    \item \textbf{Assignment Selection} --- Pick the assignment to grade from a filterable list.
    \item \textbf{Configuration} --- Set file paths, choose AI model, enter grading instructions.
    \item \textbf{Grading Progress} --- Watch live grading with a progress bar and log output.
    \item \textbf{Results} --- View scores, confidence levels, and flagged items in a sortable table.
\end{enumerate}

\begin{tipbox}
\textbf{Keyboard shortcuts:} \texttt{Tab/Shift+Tab} to navigate, \texttt{Enter} to select, \texttt{Escape} to go back, \texttt{q} to quit.
\end{tipbox}

\subsection{\texttt{gui} --- Web GUI}

\begin{lstlisting}[language=bash]
gs-autograde gui                     # Default: http://127.0.0.1:8080
gs-autograde gui --port 3000         # Custom port
\end{lstlisting}

Opens a cross-platform web interface in your browser with a 5-step grading wizard:

\begin{enumerate}[itemsep=3pt]
    \item \textbf{Login} --- Authenticate with your Gradescope credentials.
    \item \textbf{Select Assignment} --- Choose course and assignment from dropdown menus.
    \item \textbf{Configure} --- Upload question/answer PDFs, rubric YAML, set extra instructions, pick model.
    \item \textbf{Grade} --- Run grading with live progress display.
    \item \textbf{Export} --- Download results as Gradescope-compatible CSV.
\end{enumerate}

% ════════════════════════════════════════════════════════════
\section{LLM Providers}

\subsection{OpenCode Go (Default, Cloud)}

\begin{lstlisting}[style=yamlstyle]
llm:
  provider: "opencode-go"
  model: "deepseek-v4-flash"          # 1M context, fast, text-only
  multimodal_model: "mimo-v2.5"       # 1M context, text+image+video+audio
\end{lstlisting}

\begin{itemize}[itemsep=2pt]
    \item \textbf{deepseek-v4-flash}: 284B total params (13B active), 1M token context, 384K max output. Best for text-only grading.
    \item \textbf{mimo-v2.5}: 310B total params (15B active), 1M token context, 128K max output. Use when student submissions include images, diagrams, or handwriting.
\end{itemize}

Set the environment variable: \texttt{export OPENCODE\_GO\_API\_KEY="your-key"}

\subsection{LM Studio (Local, Private)}

\begin{lstlisting}[style=yamlstyle]
llm:
  provider: "lmstudio"
lmstudio:
  auto_discover: true
  base_url: "http://localhost:1234/v1"
\end{lstlisting}

\begin{enumerate}[itemsep=3pt]
    \item Install \href{https://lmstudio.ai/}{LM Studio}
    \item Download models via LM Studio's UI
    \item Start the local server (Developer $\rightarrow$ Local Server $\rightarrow$ Start)
    \item Run \texttt{gs-autograde models} to see available models
\end{enumerate}

\textbf{Recommended models for grading:}

\begin{center}
\begin{tabular}{lccp{4cm}}
\toprule
\textbf{Model} & \textbf{Size} & \textbf{Context} & \textbf{Best For} \\
\midrule
Gemma 4 12B & 12B & 131K & Balanced, general grading \\
Qwen 3.5 9B & 9B & 262K & Code/math grading \\
Gemma 4 E4B & 4.5B & 131K & Quick checks, fast iteration \\
Qwen 3.5 4B & 4B & 262K & Lightweight, resource-constrained \\
\bottomrule
\end{tabular}
\end{center}

\begin{warnbox}
\textbf{Privacy note:} When using cloud providers (OpenCode Go, OpenAI, Anthropic), student answers are sent to external servers. Check your university's data policy. LM Studio keeps all data on your machine.
\end{warnbox}

% ════════════════════════════════════════════════════════════
\section{LM Studio Auto-Management}

The LM Studio provider includes automatic server lifecycle management. No need to manually start the LM Studio desktop app before grading.

\subsection{Automatic Detection}

\begin{lstlisting}[language=python]
from gradescope_autograde.lmstudio.manager import detect_lmstudio, get_install_instructions

status = detect_lmstudio()
# Returns: {installed: bool, lms_path: str, server_running: bool, ...}
\end{lstlisting}

\subsection{Automatic Server Start}

\begin{lstlisting}[language=python]
from gradescope_autograde.lmstudio.manager import LmsManager

with LmsManager() as lm:
    lm.ensure_running()    # Starts server if not already running
    # ... use LM Studio for grading ...
# Server stops when exiting the context manager
\end{lstlisting}

\begin{tipbox}
Both the TUI and GUI use \texttt{LmsManager} automatically. When you select LM Studio as your provider, the server is started for you.
\end{tipbox}

\subsection{Installation Detection}

If LM Studio is not installed, the tool provides platform-specific instructions:

\begin{itemize}[itemsep=2pt]
    \item \textbf{macOS:} Download from \url{https://lmstudio.ai/} or run \texttt{curl -fsSL https://lmstudio.ai/install.sh | bash}
    \item \textbf{Linux:} \texttt{curl -fsSL https://lmstudio.ai/install.sh | bash}
    \item \textbf{Windows:} Download installer from \url{https://lmstudio.ai/}
\end{itemize}

% ════════════════════════════════════════════════════════════
\section{Rubric Design}

Rubrics are the most critical component of the grading pipeline. The AI's accuracy depends almost entirely on how well your rubric is structured.

\subsection{Rubric Format (YAML)}

\begin{lstlisting}[style=yamlstyle]
assignment: "Homework 1: Recursion"
total_points: 100

questions:
  - id: "q1"
    title: "Define recursion"
    max_points: 10
    type: short_answer
    rubric:
      - criterion: "Correct definition"
        points: 3
        description: "Must mention function calling itself"
      - criterion: "Base case mentioned"
        points: 3
        description: "Must mention stopping condition"
      - criterion: "Example provided"
        points: 2
      - criterion: "Clarity and terminology"
        points: 2

grading_guidelines:
  - "Award partial credit when student shows understanding"
  - "Blank → 0 points, flag for review"
  - "Off-topic → 0 points, flag for review"
\end{lstlisting}

\subsection{Best Practices}

\begin{enumerate}[itemsep=5pt]
    \item \textbf{Be specific.} ``Correct answer'' is vague. ``Must include both average O(n log n) and worst-case O(n\textsuperscript{2})'' is specific.

    \item \textbf{Use checklist-style criteria.} Research shows checklist rubrics produce more consistent AI grades than holistic rubrics.

    \item \textbf{Include point values per criterion.} The AI needs to know how to weight each aspect.

    \item \textbf{Add descriptions.} The \texttt{description} field tells the AI \textit{what to look for} in the student's answer.

    \item \textbf{Use \texttt{extra\_instructions}.} For edge cases or exceptions:
\begin{lstlisting}[style=yamlstyle]
extra_instructions: >
  Accept either O(n log n) or O(n^2) if properly justified.
  Deduct 1 point for missing base case explanation.
\end{lstlisting}

    \item \textbf{Start with a dry run.} Always test your rubric on a few submissions with \texttt{--dry-run} before running on the full class.
\end{enumerate}

% ════════════════════════════════════════════════════════════
\section{Workflow}

\subsection{Typical Grading Session}

\begin{enumerate}[itemsep=5pt]
    \item \textbf{Prepare:} Create or parse your rubric YAML.
    \item \textbf{Login:} \texttt{gs-autograde login}
    \item \textbf{Dry run:} \texttt{gs-autograde grade ... --dry-run}
    \item \textbf{Review rubric:} Check the dry-run results. Adjust rubric if needed.
    \item \textbf{Grade:} \texttt{gs-autograde grade ... --rubric rubric.yaml}
    \item \textbf{Review:} \texttt{gs-autograde review} — approve/reject flagged items.
    \item \textbf{Export:} \texttt{gs-autograde export --format gradescope}
    \item \textbf{Upload:} Either use \texttt{gs-autograde upload} or manually upload the CSV.
\end{enumerate}

\subsection{Resume After Interruption}

If grading is interrupted (network loss, crash, Ctrl+C), simply re-run the same \texttt{grade} command. The progress tracker (\texttt{.grading\_state.json}) remembers which submissions were already processed.

To force a full re-grade, delete the state file:

\begin{lstlisting}[language=bash]
rm .grading_state.json
\end{lstlisting}

\subsection{Output Structure}

\begin{lstlisting}[language=bash]
data/output/grades/
├── results.json             # Full results with breakdowns
├── gradescope_upload.csv    # For manual Gradescope upload
├── detailed_report.csv      # Per-criterion breakdown
└── review_queue.json        # Flagged submissions
\end{lstlisting}

% ════════════════════════════════════════════════════════════
\section{Architecture}

\begin{verbatim}
CLI (Click + Rich)
    │
Workflow Pipeline (orchestrator)
    │
├── Gradescope Client (HTTP scraping)
│   ├── Transport (CSRF auth, cookies, rate limiting)
│   ├── HTML Parser (course/assignment/submission extraction)
│   └── Browser Fallback (Playwright cookie extraction)
│
└── Grading Engine
    ├── LLM Providers
    │   ├── OpenCodeGoProvider (deepseek-v4-flash, mimo-v2.5)
    │   ├── LMStudioProvider (local gemma4, qwen3.5)
    │   ├── OpenAIProvider (optional)
    │   └── AnthropicProvider (optional)
    ├── Rubric Parser (YAML)
    ├── PDF Parser (pymupdf)
    └── Review Queue (confidence threshold)
\end{verbatim}

\subsection{Design Decisions}

\begin{itemize}[itemsep=3pt]
    \item \textbf{No official API.} Gradescope does not provide a public API. All interaction uses HTTP scraping of HTML pages.
    \item \textbf{CSRF authentication.} The login flow scrapes Rails CSRF tokens from the login page.
    \item \textbf{Dict-based interfaces.} Internal modules communicate via plain Python dicts, not strongly-typed models, for loose coupling.
    \item \textbf{LM Studio auto-discovery.} The LM Studio provider queries \texttt{GET /api/v1/models} in real-time to show available local models.
    \item \textbf{Confidence-based review.} Grades with confidence below the configured threshold (default 0.7) are automatically flagged for human review.
\end{itemize}

% ════════════════════════════════════════════════════════════
\section{Troubleshooting}

\subsection{Login fails}

\begin{itemize}[itemsep=3pt]
    \item Check your email/password in \texttt{config/config.yaml}
    \item Try the cookie-based login: extract \texttt{\_gradescope\_session} from your browser's DevTools $\rightarrow$ Application $\rightarrow$ Cookies
    \item Run \texttt{gs-autograde login --cookie "..."}
\end{itemize}

\subsection{Model not found}

\begin{itemize}[itemsep=3pt]
    \item Run \texttt{gs-autograde models} to see available models
    \item For LM Studio: ensure the server is running on port 1234
    \item For OpenCode Go: verify \texttt{OPENCODE\_GO\_API\_KEY} is set
\end{itemize}

\subsection{PIP install fails on macOS}

\begin{itemize}[itemsep=3pt]
    \item Ensure you're using Python 3.11+ (Python 3.14 on Homebrew may have issues)
    \item Try: \texttt{python3.12 -m venv venv \&\& source venv/bin/activate \&\& pip install -e .}
\end{itemize}

\subsection{Grading produces unexpected results}

\begin{itemize}[itemsep=3pt]
    \item Check your rubric — are the criteria specific enough?
    \item Run with \texttt{--dry-run} and review the AI's feedback
    \item Lower \texttt{review\_threshold} in config to catch more items for human review
    \item Try a different model (e.g., switch from deepseek-v4-flash to a larger local model)
\end{itemize}

\subsection{Gradescope blocks requests}

\begin{itemize}[itemsep=3pt]
    \item Increase \texttt{request\_delay} in config (try 2.0 or 3.0 seconds)
    \item Reduce \texttt{batch\_size} to process fewer submissions per run
    \item Use the browser fallback: install with \texttt{pip install -e ".[browser]"}
\end{itemize}

% ════════════════════════════════════════════════════════════
\section{FAQ}

\subsection*{Does this work with any Gradescope course?}
Yes, as long as you have instructor or TA access.

\subsection*{Is my data private?}
Cloud providers (OpenCode Go, OpenAI) receive student answers. LM Studio keeps everything on your machine.

\subsection*{How accurate is AI grading?}
Research shows LLMs match human accuracy on structured, rubric-based tasks. Always spot-check and use the review queue.

\subsection*{Can I grade multiple assignments at once?}
Run separate terminal sessions for different assignments. Each maintains its own progress state.

\subsection*{What if I need to regrade?}
Delete \texttt{.grading\_state.json} and re-run. Or edit your rubric and re-grade only flagged submissions.

% ════════════════════════════════════════════════════════════
\section{License}

Gradescope AutoGrade is licensed under the GNU General Public License v3.0 (GPL-3.0).

\vspace{1cm}
\begin{center}
\rule{0.5\textwidth}{0.5pt}

\vspace{0.5cm}
\textbf{Repository:} \url{https://github.com/tzhazuma/gradescope_autograde}

\textbf{Issues:} \url{https://github.com/tzhazuma/gradescope_autograde/issues}
\end{center}

\end{document}
